
\begin{frame}[fragile]{Au programme}

Modélisation de données, avec illustration sur des applications
analytiques
\vfill 

\begin{redbullet}
  \item Avec systèmes relationnels
  \item Avec systèmes NoSQL
  \item Problématique de très grands volumes\,: distribution et passage à l'échelle
\end{redbullet}

\vfill

Trois séances\,:


\begin{redbullet}
  \item Modélisation relationnelle et NoSQL
  \item Cassandra, modéle, interrogation, distribution
  \item ElasticSearch
\end{redbullet}

\vfill
\begin{noteblock}
On va essayer de montrer le plus possible les aspects pratiques\,!
\end{noteblock}

\end{frame}

\begin{frame}[fragile]{Séance du 4 mars 2025}

Principes comparés des modélisations relationnelle et NoSQL 
\vfill 

\begin{redbullet}
  \item La normalisation relationnelle, et Quiz
  \item L'interrogation, et mise en pratique
  \item Les bases dénormalisées, quiz et exercice pratique
  \item Application avec MongoDB, un peu de pratique
  \item (Peut-être) discussion du réplication et transactions
\end{redbullet}

\vfill

Séances suivantes\,: on parlera de Cassandra et d'Elastic Search
\end{frame}



\begin{frame}[fragile]{Ressources}

Cours généralistes au Cnam.

\begin{redbullet}
  \item Mon cours sur le modèle relationnel\,: \url{http://sql.bdpedia.fr}
  \item Mon cours sur les relationnels\,: \url{http://sys.bdpedia.fr}
  \item Mon cours sur les systèmes NoSQL\,: \url{http://b3d.bdpedia.fr}
 \end{redbullet}

\vfill

Ressources spécifiques au certificat.


\begin{redbullet}
  \item Un schéma relationnel sur une base ``Pathologies''
  \item Un échantillon de cette base (fictive)
  \item Un site d'interrogation SQL de cette base \url{http://deptfod.cnam.fr/bd/tp}
  \item Un échantillon de cette même base au format JSON (pour MongoDB, Cassandra et ElasticSearch)
  \item Un site de pratique de MongoDB en ligne.
 \end{redbullet}

\end{frame}



